% !TeX root = closed-systems-from-comparison-completeness.tex
% ============================================================
\chapter*{Reference Index of Notation}
\phantomsection
\addcontentsline{toc}{chapter}{Reference Index of Notation}
\label{ch:notation-index}
% ============================================================
This index records the recurring notation of the monograph.  It is a
reader reference; symbols with chapter-local meanings are defined where
they occur.

\begin{description}[leftmargin=3.1cm,style=nextline]
\item[\((U,\mathcal C)\)]
Comparison world; see \cref{def:intro-comparison-world}.
\item[\(U\)]
Primitive state set.
\item[\(\mathcal C\)]
Family of intrinsic binary comparison predicates.
\item[\(G\)]
Intrinsic automorphism group, usually \(G=\Aut(U,\mathcal C)\).
\item[\(X_A,X_B\)]
Factors in the comparison-complete product decomposition.
\item[\(\pi\)]
Quotient or projection map.
\item[\(\Phys\)]
Physical quotient, typically \((X_A\times X_B)/G\).
\item[\(R\)]
Coherent report.
\item[\(\mathsf B,\mathsf B_A,\mathsf B_B\)]
Boolean algebras of comparison-generated propositions.
\item[\(\UF(\mathsf B)\)]
Stone space of ultrafilters on \(\mathsf B\).
\item[\(\mathcal G(\mathcal I)\)]
Action or protocol groupoid associated with an index system.
\item[\(\Ob,\Mor\)]
Object and morphism sets.
\item[\(\Path,\Hol\)]
Path and holonomy data.
\item[\(F^k\)]
Transport filtration.
\item[\(F^2/F^3\)]
First stable visible transport-obstruction layer.
\item[\(\gr\)]
Associated graded object.
\item[\(M\)]
Realized smooth manifold.
\item[\(\nabla,\LC\)]
Connection and Levi--Civita connection in a smooth realization.
\item[\(\Riem,\Ric,\Scal\)]
Riemann curvature, Ricci tensor, and scalar curvature.
\item[\(\Jtwo\)]
Second-jet comparison map.
\item[\(\PhysH\)]
Hilbert-side physical quotient after Hilbert realization.
\item[\(\Obs\)]
Observable algebra or observable sector.
\item[\(\Chan\)]
Channel operator or channel class.
\item[\(\dimdiamond\)]
Intrinsic causal-diamond growth dimension.
\item[\textup{(SP1)--(SP5)}]
Closed-World Admissibility clauses.
\item[\textup{(T)}, \textup{(D)}, \textup{(U)}]
Torsion, detectability, and scale-coherence conditions.
\item[\textup{(R1)--(R3)}]
Additional realization data for the state-side reconstruction extension.
\end{description}

% ============================================================
\chapter*{Reference Index of Definitions and Terms}
\phantomsection
\addcontentsline{toc}{chapter}{Reference Index of Definitions and Terms}
\label{ch:definition-index}
% ============================================================
This index lists the recurring terms whose formal content controls the
argument.  For orientation-level descriptions, see
\hyperref[ch:standing-terminology]{Standing Terminology}; for theorem
use, consult the indicated chapter blocks.

\begin{description}[leftmargin=3.8cm,style=nextline]
\item[Admissibility]
Primitive and finite-to-global admissibility are developed in
\hyperref[ch:introduction]{Introduction} and \cref{ch:foundation}.
\item[Boundary data]
Scoped realization and quantitative boundary data are recorded in
\hyperref[ch:logical-status]{Logical Status and Dependency Ledger} and
\cref{app:conditional-completion}.
\item[Closed world]
The standing closure stance is fixed in
\hyperref[ch:introduction]{Introduction} and
\hyperref[ch:logical-status]{Logical Status and Dependency Ledger}.
\item[Coherent report]
Used in quotient and physical-report arguments; see
\cref{ch:rotation,ch:closure}.
\item[Comparison completeness]
Developed and derived in \cref{ch:bottom,ch:closure}.
\item[Comparison world]
Defined in \cref{def:intro-comparison-world}.
\item[Diagonal redundancy]
Developed in \cref{ch:bottom,ch:rotation}.
\item[External primitive]
Excluded at primitive level by the admissibility stance; see
\hyperref[ch:introduction]{Introduction} and
\hyperref[ch:logical-status]{Logical Status and Dependency Ledger}.
\item[Ghost sector]
Analyzed in the intrinsic--smooth interface; see \cref{ch:interface}.
\item[Intrinsic]
Used throughout for comparison-recoverable, quotient-descending, or
transport-determined structure; see
\hyperref[ch:standing-terminology]{Standing Terminology}.
\item[Local distinguishability]
Used in the derivation of rectangular completeness; see
\cref{ch:bottom}.
\item[Morphism locus]
Classified in \cref{ch:locus,ch:lift}.
\item[Object locus]
Classified in \cref{ch:locus,ch:lift}.
\item[Physical quotient]
Developed in \cref{ch:rotation,ch:closure}.
\item[Realization]
Formal status is fixed in
\hyperref[ch:logical-status]{Logical Status and Dependency Ledger};
smooth and field-theoretic realizations are developed from
\cref{ch:christoffel} onward.
\item[Rectangular completeness]
Developed in \cref{ch:bottom,ch:closure}.
\item[Regular sector]
Defined and used in the state-side reconstruction extension; see
\cref{ch:dynamics-program}.
\item[Semantic descent]
Developed in \cref{ch:rotation,ch:closure}.
\item[Transport obstruction]
Developed in \cref{ch:triangles,ch:descent_obstruction,ch:bridge}.
\item[Two-locus theorem]
Proved through \cref{ch:locus,ch:lift}.
\end{description}
